\documentclass[10pt,twocolumn]{article}
\usepackage[scaled]{helvet}
\renewcommand\familydefault{\sfdefault}
\usepackage[T1]{fontenc}
\usepackage[margin=0.7in]{geometry}
\usepackage{titlesec}
\usepackage{enumitem}
\usepackage{parskip}
\setlength{\columnsep}{0.24in}
\setlist[itemize]{leftmargin=1.2em, topsep=0.2em, itemsep=0.18em}
\titleformat{\section}{\large\bfseries}{\thesection}{1em}{}

\begin{document}
\title{\vspace{-1.6cm}AWS EC2 (Elastic Compute Cloud)}
\author{AWS ML Study Notes}
\date{}
\maketitle
\small

\section*{Overview}
EC2 gives you virtual machines on demand. You choose the instance family, operating system image, storage, networking, and scaling behavior, then AWS provides the compute capacity.

For ML teams, EC2 matters because it is the most flexible way to run custom training, inference, data processing, and supporting services when managed abstractions are too restrictive.

\section*{Core Concepts}
\begin{itemize}
\item Instance families are optimized for different workloads such as general purpose, compute optimized, memory optimized, storage optimized, or GPU accelerated.
\item An EC2 instance is usually launched from an AMI and attached to EBS volumes inside a VPC. Security groups, IAM roles, and bootstrap scripts define how it behaves after launch.
\item Pricing options include on demand, reserved, and spot. Cost and interruption tolerance should influence how you choose between them for ML workloads.
\end{itemize}

\section*{How It Works}
\begin{itemize}
\item You select an AMI, instance type, subnet, security group, IAM role, and storage layout. User data or configuration management then installs dependencies and starts services.
\item If the instance hosts a long running service, it is usually placed behind a load balancer and managed by an Auto Scaling group for replacement and horizontal growth.
\item If the instance hosts a one off training or data job, it can be provisioned, used, and terminated automatically as part of a workflow to avoid paying for idle compute.
\end{itemize}

\section*{AI and ML Relevance}
\begin{itemize}
\item EC2 is common when you need custom CUDA drivers, niche libraries, background agents, or low level operating system access that is awkward in more managed services.
\item GPU instances support deep learning training and high throughput inference, while CPU instances often handle preprocessing, orchestration, and lightweight model serving.
\end{itemize}

\section*{Operational Notes}
\begin{itemize}
\item Patch management, kernel security, disk sizing, and monitoring are your responsibility on EC2. Flexibility comes with operational ownership.
\item Use spot instances for fault tolerant training or batch processing, but combine them with checkpointing because the instance can be reclaimed.
\item Right size aggressively. ML workloads can burn money quickly when overprovisioned GPUs or memory heavy instances are left running.
\end{itemize}

\section*{What To Remember}
\begin{itemize}
\item EC2 is infrastructure, not a platform. You gain control and lose convenience.
\item EBS stores block data for the instance; S3 stores durable object data outside the instance.
\item For production systems, think in terms of replacement and automation, not logging into a single pet server.
\end{itemize}

\end{document}
